% main.tex: esqueleto do relatório técnico do Insper AI (refaz-paper e cria-paper)
% Ordem e regras das seções: padrao/estrutura.md. Redação: padrao/guia-de-estilo.md.
% Opções da classe: [english] para inglês; [numerico] para citações [1] em vez de autor-ano.
\documentclass{insperai}

\title{[Título do trabalho]}
% \subtitulo{[Subtítulo opcional]}
\author{[Autor 1] \and [Autor 2]}   % blocos centralizados; \AND quebra a linha de autores
% \afiliacao{Insper AI, Insper}      % padrão
% \email{contato@insperai.com.br}    % padrão
% \data{18 de novembro de 2026}      % padrão: hoje
\resumo{[Resumo de 150 a 250 palavras: problema, o que foi feito, resultado principal e o que ele significa.]}

\begin{document}
\maketitle

\section{Introdução}
% Contexto e problema. Termina com UM parágrafo de contribuições, sem lista.

\section{Trabalhos relacionados}

\section{Método}
% Pode ter subseções: Dados, Modelo, Protocolo, Métricas.

\section{Resultados}
% Figuras: \begin{figure}[tbp]\centering\includegraphics[width=\linewidth]{figura.pdf}
%          \legenda{Título.}{Descrição.}\label{fig:x}\end{figure}   (figure* para a largura toda)
% Tabelas: booktabs (\toprule, \midrule, \bottomrule), sem linhas verticais;
%          \linhadestaque antes da linha do método proposto.

\section{Limitações}

\section{Código e dados}
% Um parágrafo só: onde está o código (URL e commit/tag), o que tem lá, de onde vêm os dados e a licença.

\section{Conclusão}

% \section*{Contribuições dos autores}   % opcional
% \section*{Ética e impacto}             % opcional

\bibliography{refs}

% \appendix                               % opcional
% \section{Hiperparâmetros}
\end{document}
